\section*{Блок 1. Асимптотическое поведение, сходимость почти наверное и лемма Бореля-Кантелли}
\addcontentsline{toc}{section}{Блок 1. Асимптотическое поведение, сходимость почти наверное и лемма Бореля-Кантелли}

\subsection*{1. Семантический анализ и математическая генерализация}
\addcontentsline{toc}{subsection}{1. Семантический анализ и математическая генерализация}

Проанализировав опорную задачу из КР и родственные задачи семинаров, можно с уверенностью сказать: этот блок проверяет понимание тонкой разницы между видами сходимостей (по вероятности $\toP$, почти наверное $\toAS$, в $L^p$) и, самое главное, \textbf{умение работать с экстремумами независимых случайных величин}. 

Когда в задаче требуется найти с вероятностью 1 (почти наверное) предел вида $\varlimsup_{n\to\infty} f(\xi_n, n)$, стандартные законы больших чисел не работают. Здесь мы имеем дело с поведением «хвостов» распределений. Математическим ядром таких задач всегда является \textbf{Лемма Бореля-Кантелли (ЛБК)}.

Абстрактная формулировка проблемы выглядит так: дана последовательность независимых случайных величин $\xi_n$. Необходимо найти такую константу $C$, что:
\[
    \P\left(\varlimsup_{n\to\infty} \frac{\xi_n - a_n}{b_n} = C\right) = 1.
\]
Это эквивалентно поиску «критического порога», при переходе через который ряд из вероятностей $\sum \P(A_n)$ меняет свое поведение со сходящегося на расходящийся.

\subsection*{2. Фундаментальный инструментарий: Асимптотика хвостов}
\addcontentsline{toc}{subsection}{2. Фундаментальный инструментарий: Асимптотика хвостов}

Для успешного применения ЛБК вам необходимо знать, как ведут себя хвосты $\P(\xi > x)$ базовых распределений при $x \to +\infty$. Интегрировать по частям на контрольной — непозволительная трата времени. Запомните следующие эквивалентности:


\begin{itemize}
    \item \textbf{Нормальное распределение} $\N(0,1)$ (через отношение Миллса):
    \[ \P(\xi > x) = \frac{1}{\sqrt{2\pi}} \int_x^{+\infty} e^{-t^2/2} dt \sim \frac{1}{\sqrt{2\pi} \, x} e^{-x^2/2}. \]
    \item \textbf{Гамма-распределение} $\Gamma(\alpha, \lambda)$ (включая $\chi^2$ и Эрланга):
    \[ \P(\xi > x) = \frac{\lambda^\alpha}{\Gamma(\alpha)} \int_x^{+\infty} t^{\alpha-1} e^{-\lambda t} dt \sim \frac{\lambda^{\alpha-1}}{\Gamma(\alpha)} x^{\alpha-1} e^{-\lambda x}. \]
    \item \textbf{Экспоненциальное распределение} $Exp(\lambda)$ (тождественное равенство):
    \[ \P(\xi > x) = e^{-\lambda x}. \]
    \item \textbf{Распределения с тяжелыми хвостами} (Парето, Коши, Фреше). Для них убывание степенное, а не экспоненциальное:
    \[ \P(\xi > x) \sim C x^{-\alpha}. \]
    \textit{Следствие:} Для тяжелых хвостов нормирующая последовательность всегда имеет степенной вид $b_n = n^{1/\alpha}$, а не логарифмический.
\end{itemize}

\subsection*{3. Универсальный алгоритм решения}
\addcontentsline{toc}{subsection}{3. Универсальный алгоритм решения}

\begin{MethodBox}[Алгоритм поиска $\varlimsup$ через ЛБК]
Пусть требуется найти $C = \varlimsup_{n\to\infty} \eta_n$, где $\eta_n = \frac{\xi_n - a_n}{b_n}$. 
Задача решается в 4 строго выверенных шага:

\begin{enumerate}[label=\textbf{\arabic*.}]
    \item \textbf{Введение порога.} Зафиксируем произвольное $\varepsilon > 0$ и рассмотрим два типа событий:
    \[ 
        A_n^+ = \{ \eta_n > C + \varepsilon \}, \quad A_n^- = \{ \eta_n > C - \varepsilon \} 
    \]
    
    \item \textbf{Оценка вероятностей.} Выразим $\xi_n$ из неравенств и найдем асимптотику вероятностей:
    \[ 
        p_n^{\pm} = \P(\xi_n > a_n + b_n(C \pm \varepsilon)) 
    \]
    Используйте асимптотику хвостов распределений. В аргументе экспоненты появятся логарифмы (например, $a_n \sim \ln n$), которые превратят экспоненту в степень $n$.
    
    \item \textbf{Анализ рядов.} Исследуем ряды $\sum p_n^+$ и $\sum p_n^-$. Чаще всего они сводятся к обобщенному гармоническому ряду $\sum \frac{1}{n^p (\ln n)^q}$.
    
    \item \textbf{Синтез по ЛБК.} Подберите $C$ так, чтобы выполнялись два условия:
    \begin{itemize}
        \item Для $A_n^+$ (с $+\varepsilon$): ряд \textbf{сходится}. По I лемме Б-К событие произойдет конечное число раз п.н. Значит, $\varlimsup \eta_n \le C$.
        \item Для $A_n^-$ (с $-\varepsilon$): ряд \textbf{расходится} и события независимы. По II лемме Б-К событие произойдет бесконечно часто п.н. Значит, $\varlimsup \eta_n \ge C$.
    \end{itemize}
    Вывод: $\varlimsup \eta_n = C$ почти наверное.
\end{enumerate}
\end{MethodBox}

\vspace{1em}

\subsection*{4. Эвристики и неочевидные связи}
\addcontentsline{toc}{subsection}{4. Эвристики и неочевидные связи}

Разберем структуру опорной задачи из контрольной (распределение $\Gamma(2,1)$ с плотностью $p(x) = 4xe^{-2x}$). Если мы применим алгоритм, нам нужно оценить $\P(\xi_n > x)$. 
Согласно нашей асимптотике, $\P(\xi > x) \sim 2x e^{-2x}$.

Подставим границу $x_n = \frac{1}{2}\ln n + C \ln\ln n$ (временно отбросим $\varepsilon$ для прикидки):
\begin{align*}
    \P(\xi_n > x_n) &\sim 2 \left( \frac{1}{2}\ln n + C\ln\ln n \right) \exp\left\{ -2 \left( \frac{1}{2}\ln n + C \ln\ln n \right) \right\} \\
    &\sim (\ln n) \cdot \exp\left\{ -\ln n - 2C\ln\ln n \right\} \\
    &= \ln n \cdot \frac{1}{n} \cdot \frac{1}{(\ln n)^{2C}} = \frac{1}{n (\ln n)^{2C-1}}.
\end{align*}

\subsubsection*{Обобщенный признак Бертрана (Логарифмическая шкала)}
При подстановке нормирующих последовательностей вида $x_n = a \ln n + b \ln \ln n$ в экспоненциальные хвосты, вероятности $\P(\xi > x_n)$ трансформируются в члены обобщенного гармонического ряда. 

Анализ таких рядов опирается на \textbf{шкалу логарифмических классов} (интегральный признак Коши-Маклорена / обобщенный признак Бертрана).
Рассмотрим ряд вида:
\[ \sum_{n=N}^{\infty} \frac{1}{n (\ln n) (\ln \ln n) \dots (\ln^{(k-1)} n) (\ln^{(k)} n)^p}, \]
где $\ln^{(k)} n$ обозначает $k$-ю итерацию логарифма. 
Данный ряд \textbf{сходится} тогда и только тогда, когда показатель степени при последнем сомножителе $p > 1$, и \textbf{расходится} при $p \le 1$.

\begin{HackBox}[Абсолютная формула критического порога (Мастер-теорема)]
Любая задача на поиск константы $C$ в пределе $\varlimsup_{n\to\infty} \frac{\xi_n - a_n}{b_n} = C$ для распределений экспоненциального типа сводится к строгой алгебраической формуле, минуя анализ рядов.

Пусть асимптотика хвоста имеет вид:
\[ \P(\xi_n > x) \sim K \cdot x^\beta e^{-\lambda x}. \]

Для нормировки $a_n = \frac{1}{\lambda}\ln n$ и $b_n = \frac{1}{\lambda} \ln \ln n$, подстановка $x_n = a_n + C \cdot b_n$ генерирует общий член ряда:
\[ \P(\xi_n > x_n) \propto \frac{(\ln n)^\beta}{n (\ln n)^C} = \frac{1}{n (\ln n)^{C - \beta}}. \]

Согласно обобщенному признаку Бертрана, критическая точка (граница между сходимостью и расходимостью, определяющая верхний предел) всегда достигается при равенстве степени знаменателя единице:
\[ C - \beta = 1 \implies C = \beta + 1. \]

\textbf{Сводка быстрых ответов (при стандартной нормировке $\frac{\xi_n - \frac{1}{\lambda}\ln n}{\frac{1}{\lambda}\ln\ln n}$):}
\begin{enumerate}
    \item Для $Exp(\lambda)$: $\beta = 0 \implies C = 1$.
    \item Для $\Gamma(\alpha, \lambda)$: $\beta = \alpha - 1 \implies C = \alpha$.
    \item Для $\N(0,1)$ (при нормировке $a_n = \sqrt{2\ln n}$): показатель перед экспонентой равен $x^{-1}$, следовательно $\beta = -1$, откуда возникает классическая поправка $C = -\frac{1}{2}$ (для логарифма в знаменателе).
\end{enumerate}
Использование данного инварианта позволяет вычислять $\varlimsup$ мгновенно, проверяя лишь степенной множитель перед экспонентой в асимптотическом разложении плотности.

\textbf{Универсальная теорема для экспоненциальных хвостов:}
Если $\P(\xi_n > x) \sim K x^\beta e^{-\lambda x}$, то для нахождения предела вида $\varlimsup \frac{\xi_n - \frac{1}{\lambda}\ln n}{\ln \ln n}$ значение константы $C$ всегда равно $\frac{\beta}{\lambda}$.
В нашей задаче $\beta = 1, \lambda = 2$, откуда мгновенно следует ответ $C = 1/2$ (однако в условии нормализация $\frac{1}{2}\ln n$, поэтому при $b_n = \ln\ln n$ слагаемое $2C$ балансирует $\beta=1$, и ответ $C=1$). Строгий контроль коэффициентов исключает вычислительные ошибки.
\end{HackBox}

\subsection*{5. Опасные места и контрпримеры (Теоретическая часть)}
\addcontentsline{toc}{subsection}{5. Опасные места и контрпримеры (Теоретическая часть)}

Помимо счетных задач, Блок 1 содержит мощный пласт теоретических контрпримеров (задачи 1, 3, 8). Это маркер того, что на контрольной могут быть теоретические вопросы на понимание.

\begin{WarningBox}[Сходимость по вероятности $\not\implies$ Сходимость почти наверное]
Студенты часто пытаются использовать свойства сходимости $\toP$ там, где требуется $\toAS$. Запомните классический \textbf{контрпример «бегущего прямоугольника»}:
Пусть отрезок $[0, 1]$ разбит на интервалы длины $1/2, 1/3, 1/4$ и т.д., которые «бегают» по кругу. Индикатор попадания в такой интервал $\xi_n \toP 0$, так как длина интервала стремится к 0. Но $\xi_n \not\toAS 0$, потому что любая точка $x \in [0,1]$ покрывается интервалами бесконечно часто (ряд длин расходится: $\sum 1/n = \infty$, II лемма Б-К). 

\textit{Секрет:} $\toP \implies \toAS$ работает \textbf{только} если ряд вероятностей отклонений сходится (это напрямую следует из I леммы Б-К). В дискретных пространствах сходимости совпадают, так как вероятность не может быть сколь угодно малой, не став нулем.
\end{WarningBox}
